\section{Related Works} \label{sec:related}
\subsection{Talking Avatar Generation}
Talking avatar generation aims to generate a lifelike talking avatar video from a given reference image and an audio. Earlier methods in this field had focused on synthesizing accurate lip movement from the driving audio~\citep{wav2lip, makeittalk, stylelipsync, stylesync, keysync}. These approaches are extended to generate holistic head movements, including a rhythmical head motion~\citep{makeittalk, emo, float, vasa_1} and vivid facial expressions, such as eye blink and jaw movement~\citep{sadtalker, megaportrait, emoportraits, vasa_1, liveportrait}. SadTalker~\citep{sadtalker}, for instance, leverages 3D morphable models (3DMM)~\citep{bfm} as an intermediate representation for the non-verbal facial expression. EMO~\citep{emo} and its subsequent models~\citep{loopy, hallo3, echomimic} leverage the foundation image diffusion model (e.g., StableDiffusion~\citep{ldm}) for photorealistic portrait animation. 
Recent works~\citep{vasa_1, anitalker, float} introduce generative modeling techniques, such as diffusion models or flow matching into a learned motion space~\citep{megaportrait, lia}, achieving real-time head motion generation.

\subsection{Listening Avatar Generation}
Another line of works~\citep{l2l, rlhg, mfr_net, pch, customlistener, ditailistener} focus on generating realistic listening head motions, such as nodding or focusing. Generating such responsive motions is challenging because the relationship between the speaker and listener is inherently one-to-many~\citep{l2l, diffusion_listening}, and the cues are context-dependent and weakly supervised.
Hence, most of the works generate the \textit{personalized} listening motion~\citep{l2l, rlhg} or leverage explicit control signals, such as \textit{text instruction}~\citep{customlistener} and \textit{pose prior}~\citep{diffusion_listening}. 

\subsection{Dyadic Conversational Avatar Generation}
Recently, several studies investigated dyadic motion generation~\citep{dim, infp, arig}, which involves modeling the interactive behavior of two participants engaged in a conversational setting.
DIM~\citep{dim} quantizes dyadic motions into two discrete latent spaces~\citep{vqvae}, one for the verbal and one for the non-verbal. However, it requires a manual role-switching signal between the two spaces, resulting in discontinuous transitions between speaker and listener states. INFP~\citep{infp} addresses this limitation by introducing audio-conditioned verbal and non-verbal memory banks within a unified motion generator. 
However, this is ill-suited for real-time interactive generation as its bidirectional transformer~\citep{dit} requires access to the entire context of the conversation.
ARIG~\citep{arig} generates implicit 3D keypoints~\citep{liveportrait} as its motion representation, primarily focusing on facial expressions. Yet it struggles with temporal consistency and fails to produce holistic head motions.

In this work, we focus on modeling the interactive behavior of both participants continuously influencing each other through verbal and non-verbal signals. We generate holistic interactive motions, for instance, talking, head movement, listening, and focusing, using a diffusion forcing framework that promptly responds to the multimodal user signals. 
